% json2dir.tex -- the json2dir evaluator, written for LuaTeX (RFC J2D-1).
%
% Run by the `json2dir` launcher in a sandbox (--safer, no shell escape):
%   luatex --ini --safer --no-shell-escape -interaction=batchmode -output-directory=DIR \
%     '\catcode123=1 \catcode125=2 \def\inputfile{input.json}% json2dir.tex -- the json2dir evaluator, written for LuaTeX (RFC J2D-1).
%
% Run by the `json2dir` launcher in a sandbox (--safer, no shell escape):
%   luatex --ini --safer --no-shell-escape -interaction=batchmode -output-directory=DIR \
%     '\catcode123=1 \catcode125=2 \def\inputfile{input.json}% json2dir.tex -- the json2dir evaluator, written for LuaTeX (RFC J2D-1).
%
% Run by the `json2dir` launcher in a sandbox (--safer, no shell escape):
%   luatex --ini --safer --no-shell-escape -interaction=batchmode -output-directory=DIR \
%     '\catcode123=1 \catcode125=2 \def\inputfile{input.json}% json2dir.tex -- the json2dir evaluator, written for LuaTeX (RFC J2D-1).
%
% Run by the `json2dir` launcher in a sandbox (--safer, no shell escape):
%   luatex --ini --safer --no-shell-escape -interaction=batchmode -output-directory=DIR \
%     '\catcode123=1 \catcode125=2 \def\inputfile{input.json}\input{json2dir.tex}'
% It reads \inputfile, validates the whole document and writes DIR/json2dir-tree.sh, a
% POSIX shell script that creates the tree. Any error: message on stderr, exit 1, no script.
%
% TeX drives; the JSON work is done by LuaTeX's embedded Lua in \directlua. --safer
% disables io.open (it uses io.input to read), so the script is written with \openout/\write: LuaTeX
% writes characters to \write files as raw UTF-8 bytes (no ^^ notation, no line wrapping),
% \newlinechar=10 turns LF into LF, and tex.sprint(-2, ...) turns every byte into a plain
% "other" character token, so the bytes come out exactly as generated.
\catcode`\{=1 \catcode`\}=2
%
% Read the Lua code with @ as the escape character and \ % and LF as ordinary characters,
% so backslashes, Lua patterns and "--" comments survive. Braces must stay balanced.
\begingroup \catcode`\%=12 \catcode10=12 \endlinechar=10 \catcode`\@=0 \catcode`\\=12
@directlua{
local byte, sub, char, find = string.byte, string.sub, string.char, string.find

local function fail(msg) error({ msg = msg }, 0) end

-- RFC 3: strict UTF-8
local function check_utf8(s)
  local i, n = 1, #s
  while i <= n do
    local c = byte(s, i)
    if c < 0x80 then
      i = i + 1
    else
      local len, min
      if c >= 0xC2 and c <= 0xDF then len, min = 2, 0x80
      elseif c >= 0xE0 and c <= 0xEF then len, min = 3, 0x800
      elseif c >= 0xF0 and c <= 0xF4 then len, min = 4, 0x10000
      else fail("input is not valid UTF-8") end
      if i + len - 1 > n then fail("input is not valid UTF-8") end
      local cp = c & (0x7F >> len)
      for k = 1, len - 1 do
        local d = byte(s, i + k)
        if d < 0x80 or d > 0xBF then fail("input is not valid UTF-8") end
        cp = cp * 64 + (d - 0x80)
      end
      if cp < min or cp > 0x10FFFF or (cp >= 0xD800 and cp <= 0xDFFF) then
        fail("input is not valid UTF-8")
      end
      i = i + len
    end
  end
end

-- RFC 3, 4: strict JSON. Objects -> {obj=true, keys=..., map=...}, arrays -> {arr=true, ...},
-- strings -> Lua strings. Numbers, true, false, null are rejected right away (RFC 4.5).
local function parse(s)
  local pos = 1
  local escapes = { [34] = '"', [92] = "\\", [47] = "/", [98] = "\b", [102] = "\f",
    [110] = "\n", [114] = "\r", [116] = "\t" }

  local function skip_ws() pos = find(s, "[^ \t\r\n]", pos) or #s + 1 end
  local function bad(what) fail("input is not valid JSON: " .. what .. " at byte " .. pos) end

  local function hex4()
    local h = sub(s, pos, pos + 3)
    if not h:match("^%x%x%x%x$") then bad("bad \\u escape") end
    pos = pos + 4
    return tonumber(h, 16)
  end

  local function parse_string()
    pos = pos + 1
    local buf = {}
    while true do
      local j = find(s, '["\\%c]', pos)
      if not j then bad("unterminated string") end
      if j > pos then buf[#buf + 1] = sub(s, pos, j - 1) end
      local c = byte(s, j)
      pos = j + 1
      if c == 34 then
        return table.concat(buf)
      elseif c == 92 then
        local e = byte(s, pos)
        if e == 117 then
          pos = pos + 1
          local cp = hex4()
          if cp >= 0xD800 and cp <= 0xDBFF then
            if sub(s, pos, pos + 1) ~= "\\u" then bad("unpaired surrogate") end
            pos = pos + 2
            local lo = hex4()
            if lo < 0xDC00 or lo > 0xDFFF then bad("unpaired surrogate") end
            cp = 0x10000 + (cp - 0xD800) * 0x400 + (lo - 0xDC00)
          elseif cp >= 0xDC00 and cp <= 0xDFFF then
            bad("unpaired surrogate")
          end
          buf[#buf + 1] = utf8.char(cp)
        elseif e and escapes[e] then
          buf[#buf + 1] = escapes[e]
          pos = pos + 1
        else
          bad("bad escape")
        end
      elseif c == 127 then
        buf[#buf + 1] = "\127"
      else
        bad("control character in string")
      end
    end
  end

  local parse_value

  local function parse_object(depth)
    pos = pos + 1
    local o = { obj = true, keys = {}, map = {} }
    skip_ws()
    if byte(s, pos) == 125 then pos = pos + 1; return o end
    while true do
      skip_ws()
      if byte(s, pos) ~= 34 then bad("expected member name") end
      local name = parse_string()
      skip_ws()
      if byte(s, pos) ~= 58 then bad("expected ':'") end
      pos = pos + 1
      local v = parse_value(depth + 1)
      if o.map[name] == nil then o.keys[#o.keys + 1] = name end
      o.map[name] = v -- duplicate names: the last occurrence wins (RFC 4.2)
      skip_ws()
      local c = byte(s, pos)
      pos = pos + 1
      if c == 125 then return o end
      if c ~= 44 then pos = pos - 1; bad("expected ',' or end of object") end
    end
  end

  local function parse_array(depth)
    pos = pos + 1
    local a = { arr = true }
    skip_ws()
    if byte(s, pos) == 93 then pos = pos + 1; return a end
    while true do
      a[#a + 1] = parse_value(depth + 1)
      skip_ws()
      local c = byte(s, pos)
      pos = pos + 1
      if c == 93 then return a end
      if c ~= 44 then pos = pos - 1; bad("expected ',' or ']'") end
    end
  end

  function parse_value(depth)
    if depth > 1000 then bad("nesting too deep") end
    skip_ws()
    local c = byte(s, pos)
    if c == 123 then return parse_object(depth)
    elseif c == 91 then return parse_array(depth)
    elseif c == 34 then return parse_string()
    elseif c == nil then bad("unexpected end of input")
    elseif find(s, "^[-0-9tfn]", pos) then bad("numbers, true, false and null are not allowed")
    else bad("unexpected character") end
  end

  if sub(s, 1, 3) == "\239\187\191" then pos = 4 end -- optional BOM
  local v = parse_value(0)
  skip_ws()
  if pos <= #s then bad("trailing data") end
  return v
end

-- names in ascending order of their UTF-8 bytes, whatever the locale
local function byteless(a, b)
  for i = 1, math.min(#a, #b) do
    local x, y = byte(a, i), byte(b, i)
    if x ~= y then return x < y end
  end
  return #a < #b
end

local function quote(s) return "'" .. (s:gsub("'", "'\\''")) .. "'" end

-- RFC 4 (validation of the whole document) and RFC 5 (as shell code), in one walk.
local function walk(o, dir, out)
  local keys = {}
  for _, k in ipairs(o.keys) do keys[#keys + 1] = k end
  table.sort(keys, byteless)
  for _, name in ipairs(keys) do
    local path = dir .. "/" .. name
    if name == "" or name == "." or name == ".." or find(name, "/", 1, true) or find(name, "\0", 1, true) then
      fail(path .. ": invalid name")
    end
    local v, q = o.map[name], quote(path)
    out[#out + 1] = "_json2dir_remove " .. q
    if type(v) == "string" then
      if find(v, "\0", 1, true) then fail(path .. ": content with NUL cannot be passed to the shell") end
      out[#out + 1] = "_json2dir_file " .. q .. " " .. quote(v)
    elseif v.arr then
      if #v ~= 2 or type(v[1]) ~= "string" or type(v[2]) ~= "string" then
        fail(path .. ': an array must be ["link", target] or ["script", content]')
      end
      if v[1] ~= "link" and v[1] ~= "script" then fail(path .. ": unknown array kind " .. v[1]) end
      if find(v[2], "\0", 1, true) then fail(path .. ": " .. v[1] .. " payload with NUL cannot be passed to the shell") end
      out[#out + 1] = "_json2dir_" .. v[1] .. " " .. q .. " " .. quote(v[2])
    else
      out[#out + 1] = "_json2dir_directory " .. q
      walk(v, path, out)
    end
  end
end

-- Only constant code here; every name, content and target is a single-quoted word.
local helpers = [[
set -eu
_json2dir_error() {
  printf 'Error: couldn\047t create %s at %s.\n' "$1" "$2" >&2
  exit 1
}
_json2dir_remove() {
  rm -f -- "$1" 2>/dev/null || :
}
_json2dir_directory() {
  if [ ! -d "$1" ]; then
    mkdir -- "$1" || _json2dir_error directory "$1"
  fi
  (cd -- "$1") || _json2dir_error directory "$1"
}
_json2dir_file() {
  printf '%s' "$2" > "$1" || _json2dir_error file "$1"
}
_json2dir_script() {
  printf '%s' "$2" > "$1" || _json2dir_error script "$1"
  chmod a+x "$1" || _json2dir_error script "$1"
}
_json2dir_link() {
  if [ -d "$1" ] && [ ! -L "$1" ]; then
    _json2dir_error symlink "$1"
  fi
  ln -s -- "$2" "$1" || _json2dir_error symlink "$1"
}]]

json2dir = {}

-- Called from TeX: read, check and translate the input; on any error report and exit 1.
function json2dir.run(file)
  local ok, err = pcall(function()
    -- --safer disables io.open entirely, but io.input(name) still opens a file for
    -- reading (fopen "r": on POSIX the bytes come through unchanged).
    local okf, f = pcall(io.input, file)
    if not okf or not f then fail("couldn't read the input file") end
    local s = f:read("a")
    f:close()
    if not s then fail("couldn't read the input file") end
    check_utf8(s)
    local root = parse(s)
    if type(root) ~= "table" or not root.obj then fail("the root must be an object") end
    local out = { helpers }
    walk(root, ".", out)
    json2dir.script = table.concat(out, "\n") -- \write adds the final newline
  end)
  if not ok then
    io.stderr:write("json2dir: " .. (type(err) == "table" and err.msg or tostring(err)) .. "\n")
    io.stderr:flush()
    os.exit(1)
  end
end
}@endgroup
%
\newlinechar=10
\directlua{json2dir.run(token.get_macro("inputfile"))}
\immediate\openout1=json2dir-tree.sh
\immediate\write1{\directlua{tex.sprint(-2, json2dir.script)}}
\immediate\closeout1
\end
'
% It reads \inputfile, validates the whole document and writes DIR/json2dir-tree.sh, a
% POSIX shell script that creates the tree. Any error: message on stderr, exit 1, no script.
%
% TeX drives; the JSON work is done by LuaTeX's embedded Lua in \directlua. --safer
% disables io.open (it uses io.input to read), so the script is written with \openout/\write: LuaTeX
% writes characters to \write files as raw UTF-8 bytes (no ^^ notation, no line wrapping),
% \newlinechar=10 turns LF into LF, and tex.sprint(-2, ...) turns every byte into a plain
% "other" character token, so the bytes come out exactly as generated.
\catcode`\{=1 \catcode`\}=2
%
% Read the Lua code with @ as the escape character and \ % and LF as ordinary characters,
% so backslashes, Lua patterns and "--" comments survive. Braces must stay balanced.
\begingroup \catcode`\%=12 \catcode10=12 \endlinechar=10 \catcode`\@=0 \catcode`\\=12
@directlua{
local byte, sub, char, find = string.byte, string.sub, string.char, string.find

local function fail(msg) error({ msg = msg }, 0) end

-- RFC 3: strict UTF-8
local function check_utf8(s)
  local i, n = 1, #s
  while i <= n do
    local c = byte(s, i)
    if c < 0x80 then
      i = i + 1
    else
      local len, min
      if c >= 0xC2 and c <= 0xDF then len, min = 2, 0x80
      elseif c >= 0xE0 and c <= 0xEF then len, min = 3, 0x800
      elseif c >= 0xF0 and c <= 0xF4 then len, min = 4, 0x10000
      else fail("input is not valid UTF-8") end
      if i + len - 1 > n then fail("input is not valid UTF-8") end
      local cp = c & (0x7F >> len)
      for k = 1, len - 1 do
        local d = byte(s, i + k)
        if d < 0x80 or d > 0xBF then fail("input is not valid UTF-8") end
        cp = cp * 64 + (d - 0x80)
      end
      if cp < min or cp > 0x10FFFF or (cp >= 0xD800 and cp <= 0xDFFF) then
        fail("input is not valid UTF-8")
      end
      i = i + len
    end
  end
end

-- RFC 3, 4: strict JSON. Objects -> {obj=true, keys=..., map=...}, arrays -> {arr=true, ...},
-- strings -> Lua strings. Numbers, true, false, null are rejected right away (RFC 4.5).
local function parse(s)
  local pos = 1
  local escapes = { [34] = '"', [92] = "\\", [47] = "/", [98] = "\b", [102] = "\f",
    [110] = "\n", [114] = "\r", [116] = "\t" }

  local function skip_ws() pos = find(s, "[^ \t\r\n]", pos) or #s + 1 end
  local function bad(what) fail("input is not valid JSON: " .. what .. " at byte " .. pos) end

  local function hex4()
    local h = sub(s, pos, pos + 3)
    if not h:match("^%x%x%x%x$") then bad("bad \\u escape") end
    pos = pos + 4
    return tonumber(h, 16)
  end

  local function parse_string()
    pos = pos + 1
    local buf = {}
    while true do
      local j = find(s, '["\\%c]', pos)
      if not j then bad("unterminated string") end
      if j > pos then buf[#buf + 1] = sub(s, pos, j - 1) end
      local c = byte(s, j)
      pos = j + 1
      if c == 34 then
        return table.concat(buf)
      elseif c == 92 then
        local e = byte(s, pos)
        if e == 117 then
          pos = pos + 1
          local cp = hex4()
          if cp >= 0xD800 and cp <= 0xDBFF then
            if sub(s, pos, pos + 1) ~= "\\u" then bad("unpaired surrogate") end
            pos = pos + 2
            local lo = hex4()
            if lo < 0xDC00 or lo > 0xDFFF then bad("unpaired surrogate") end
            cp = 0x10000 + (cp - 0xD800) * 0x400 + (lo - 0xDC00)
          elseif cp >= 0xDC00 and cp <= 0xDFFF then
            bad("unpaired surrogate")
          end
          buf[#buf + 1] = utf8.char(cp)
        elseif e and escapes[e] then
          buf[#buf + 1] = escapes[e]
          pos = pos + 1
        else
          bad("bad escape")
        end
      elseif c == 127 then
        buf[#buf + 1] = "\127"
      else
        bad("control character in string")
      end
    end
  end

  local parse_value

  local function parse_object(depth)
    pos = pos + 1
    local o = { obj = true, keys = {}, map = {} }
    skip_ws()
    if byte(s, pos) == 125 then pos = pos + 1; return o end
    while true do
      skip_ws()
      if byte(s, pos) ~= 34 then bad("expected member name") end
      local name = parse_string()
      skip_ws()
      if byte(s, pos) ~= 58 then bad("expected ':'") end
      pos = pos + 1
      local v = parse_value(depth + 1)
      if o.map[name] == nil then o.keys[#o.keys + 1] = name end
      o.map[name] = v -- duplicate names: the last occurrence wins (RFC 4.2)
      skip_ws()
      local c = byte(s, pos)
      pos = pos + 1
      if c == 125 then return o end
      if c ~= 44 then pos = pos - 1; bad("expected ',' or end of object") end
    end
  end

  local function parse_array(depth)
    pos = pos + 1
    local a = { arr = true }
    skip_ws()
    if byte(s, pos) == 93 then pos = pos + 1; return a end
    while true do
      a[#a + 1] = parse_value(depth + 1)
      skip_ws()
      local c = byte(s, pos)
      pos = pos + 1
      if c == 93 then return a end
      if c ~= 44 then pos = pos - 1; bad("expected ',' or ']'") end
    end
  end

  function parse_value(depth)
    if depth > 1000 then bad("nesting too deep") end
    skip_ws()
    local c = byte(s, pos)
    if c == 123 then return parse_object(depth)
    elseif c == 91 then return parse_array(depth)
    elseif c == 34 then return parse_string()
    elseif c == nil then bad("unexpected end of input")
    elseif find(s, "^[-0-9tfn]", pos) then bad("numbers, true, false and null are not allowed")
    else bad("unexpected character") end
  end

  if sub(s, 1, 3) == "\239\187\191" then pos = 4 end -- optional BOM
  local v = parse_value(0)
  skip_ws()
  if pos <= #s then bad("trailing data") end
  return v
end

-- names in ascending order of their UTF-8 bytes, whatever the locale
local function byteless(a, b)
  for i = 1, math.min(#a, #b) do
    local x, y = byte(a, i), byte(b, i)
    if x ~= y then return x < y end
  end
  return #a < #b
end

local function quote(s) return "'" .. (s:gsub("'", "'\\''")) .. "'" end

-- RFC 4 (validation of the whole document) and RFC 5 (as shell code), in one walk.
local function walk(o, dir, out)
  local keys = {}
  for _, k in ipairs(o.keys) do keys[#keys + 1] = k end
  table.sort(keys, byteless)
  for _, name in ipairs(keys) do
    local path = dir .. "/" .. name
    if name == "" or name == "." or name == ".." or find(name, "/", 1, true) or find(name, "\0", 1, true) then
      fail(path .. ": invalid name")
    end
    local v, q = o.map[name], quote(path)
    out[#out + 1] = "_json2dir_remove " .. q
    if type(v) == "string" then
      if find(v, "\0", 1, true) then fail(path .. ": content with NUL cannot be passed to the shell") end
      out[#out + 1] = "_json2dir_file " .. q .. " " .. quote(v)
    elseif v.arr then
      if #v ~= 2 or type(v[1]) ~= "string" or type(v[2]) ~= "string" then
        fail(path .. ': an array must be ["link", target] or ["script", content]')
      end
      if v[1] ~= "link" and v[1] ~= "script" then fail(path .. ": unknown array kind " .. v[1]) end
      if find(v[2], "\0", 1, true) then fail(path .. ": " .. v[1] .. " payload with NUL cannot be passed to the shell") end
      out[#out + 1] = "_json2dir_" .. v[1] .. " " .. q .. " " .. quote(v[2])
    else
      out[#out + 1] = "_json2dir_directory " .. q
      walk(v, path, out)
    end
  end
end

-- Only constant code here; every name, content and target is a single-quoted word.
local helpers = [[
set -eu
_json2dir_error() {
  printf 'Error: couldn\047t create %s at %s.\n' "$1" "$2" >&2
  exit 1
}
_json2dir_remove() {
  rm -f -- "$1" 2>/dev/null || :
}
_json2dir_directory() {
  if [ ! -d "$1" ]; then
    mkdir -- "$1" || _json2dir_error directory "$1"
  fi
  (cd -- "$1") || _json2dir_error directory "$1"
}
_json2dir_file() {
  printf '%s' "$2" > "$1" || _json2dir_error file "$1"
}
_json2dir_script() {
  printf '%s' "$2" > "$1" || _json2dir_error script "$1"
  chmod a+x "$1" || _json2dir_error script "$1"
}
_json2dir_link() {
  if [ -d "$1" ] && [ ! -L "$1" ]; then
    _json2dir_error symlink "$1"
  fi
  ln -s -- "$2" "$1" || _json2dir_error symlink "$1"
}]]

json2dir = {}

-- Called from TeX: read, check and translate the input; on any error report and exit 1.
function json2dir.run(file)
  local ok, err = pcall(function()
    -- --safer disables io.open entirely, but io.input(name) still opens a file for
    -- reading (fopen "r": on POSIX the bytes come through unchanged).
    local okf, f = pcall(io.input, file)
    if not okf or not f then fail("couldn't read the input file") end
    local s = f:read("a")
    f:close()
    if not s then fail("couldn't read the input file") end
    check_utf8(s)
    local root = parse(s)
    if type(root) ~= "table" or not root.obj then fail("the root must be an object") end
    local out = { helpers }
    walk(root, ".", out)
    json2dir.script = table.concat(out, "\n") -- \write adds the final newline
  end)
  if not ok then
    io.stderr:write("json2dir: " .. (type(err) == "table" and err.msg or tostring(err)) .. "\n")
    io.stderr:flush()
    os.exit(1)
  end
end
}@endgroup
%
\newlinechar=10
\directlua{json2dir.run(token.get_macro("inputfile"))}
\immediate\openout1=json2dir-tree.sh
\immediate\write1{\directlua{tex.sprint(-2, json2dir.script)}}
\immediate\closeout1
\end
'
% It reads \inputfile, validates the whole document and writes DIR/json2dir-tree.sh, a
% POSIX shell script that creates the tree. Any error: message on stderr, exit 1, no script.
%
% TeX drives; the JSON work is done by LuaTeX's embedded Lua in \directlua. --safer
% disables io.open (it uses io.input to read), so the script is written with \openout/\write: LuaTeX
% writes characters to \write files as raw UTF-8 bytes (no ^^ notation, no line wrapping),
% \newlinechar=10 turns LF into LF, and tex.sprint(-2, ...) turns every byte into a plain
% "other" character token, so the bytes come out exactly as generated.
\catcode`\{=1 \catcode`\}=2
%
% Read the Lua code with @ as the escape character and \ % and LF as ordinary characters,
% so backslashes, Lua patterns and "--" comments survive. Braces must stay balanced.
\begingroup \catcode`\%=12 \catcode10=12 \endlinechar=10 \catcode`\@=0 \catcode`\\=12
@directlua{
local byte, sub, char, find = string.byte, string.sub, string.char, string.find

local function fail(msg) error({ msg = msg }, 0) end

-- RFC 3: strict UTF-8
local function check_utf8(s)
  local i, n = 1, #s
  while i <= n do
    local c = byte(s, i)
    if c < 0x80 then
      i = i + 1
    else
      local len, min
      if c >= 0xC2 and c <= 0xDF then len, min = 2, 0x80
      elseif c >= 0xE0 and c <= 0xEF then len, min = 3, 0x800
      elseif c >= 0xF0 and c <= 0xF4 then len, min = 4, 0x10000
      else fail("input is not valid UTF-8") end
      if i + len - 1 > n then fail("input is not valid UTF-8") end
      local cp = c & (0x7F >> len)
      for k = 1, len - 1 do
        local d = byte(s, i + k)
        if d < 0x80 or d > 0xBF then fail("input is not valid UTF-8") end
        cp = cp * 64 + (d - 0x80)
      end
      if cp < min or cp > 0x10FFFF or (cp >= 0xD800 and cp <= 0xDFFF) then
        fail("input is not valid UTF-8")
      end
      i = i + len
    end
  end
end

-- RFC 3, 4: strict JSON. Objects -> {obj=true, keys=..., map=...}, arrays -> {arr=true, ...},
-- strings -> Lua strings. Numbers, true, false, null are rejected right away (RFC 4.5).
local function parse(s)
  local pos = 1
  local escapes = { [34] = '"', [92] = "\\", [47] = "/", [98] = "\b", [102] = "\f",
    [110] = "\n", [114] = "\r", [116] = "\t" }

  local function skip_ws() pos = find(s, "[^ \t\r\n]", pos) or #s + 1 end
  local function bad(what) fail("input is not valid JSON: " .. what .. " at byte " .. pos) end

  local function hex4()
    local h = sub(s, pos, pos + 3)
    if not h:match("^%x%x%x%x$") then bad("bad \\u escape") end
    pos = pos + 4
    return tonumber(h, 16)
  end

  local function parse_string()
    pos = pos + 1
    local buf = {}
    while true do
      local j = find(s, '["\\%c]', pos)
      if not j then bad("unterminated string") end
      if j > pos then buf[#buf + 1] = sub(s, pos, j - 1) end
      local c = byte(s, j)
      pos = j + 1
      if c == 34 then
        return table.concat(buf)
      elseif c == 92 then
        local e = byte(s, pos)
        if e == 117 then
          pos = pos + 1
          local cp = hex4()
          if cp >= 0xD800 and cp <= 0xDBFF then
            if sub(s, pos, pos + 1) ~= "\\u" then bad("unpaired surrogate") end
            pos = pos + 2
            local lo = hex4()
            if lo < 0xDC00 or lo > 0xDFFF then bad("unpaired surrogate") end
            cp = 0x10000 + (cp - 0xD800) * 0x400 + (lo - 0xDC00)
          elseif cp >= 0xDC00 and cp <= 0xDFFF then
            bad("unpaired surrogate")
          end
          buf[#buf + 1] = utf8.char(cp)
        elseif e and escapes[e] then
          buf[#buf + 1] = escapes[e]
          pos = pos + 1
        else
          bad("bad escape")
        end
      elseif c == 127 then
        buf[#buf + 1] = "\127"
      else
        bad("control character in string")
      end
    end
  end

  local parse_value

  local function parse_object(depth)
    pos = pos + 1
    local o = { obj = true, keys = {}, map = {} }
    skip_ws()
    if byte(s, pos) == 125 then pos = pos + 1; return o end
    while true do
      skip_ws()
      if byte(s, pos) ~= 34 then bad("expected member name") end
      local name = parse_string()
      skip_ws()
      if byte(s, pos) ~= 58 then bad("expected ':'") end
      pos = pos + 1
      local v = parse_value(depth + 1)
      if o.map[name] == nil then o.keys[#o.keys + 1] = name end
      o.map[name] = v -- duplicate names: the last occurrence wins (RFC 4.2)
      skip_ws()
      local c = byte(s, pos)
      pos = pos + 1
      if c == 125 then return o end
      if c ~= 44 then pos = pos - 1; bad("expected ',' or end of object") end
    end
  end

  local function parse_array(depth)
    pos = pos + 1
    local a = { arr = true }
    skip_ws()
    if byte(s, pos) == 93 then pos = pos + 1; return a end
    while true do
      a[#a + 1] = parse_value(depth + 1)
      skip_ws()
      local c = byte(s, pos)
      pos = pos + 1
      if c == 93 then return a end
      if c ~= 44 then pos = pos - 1; bad("expected ',' or ']'") end
    end
  end

  function parse_value(depth)
    if depth > 1000 then bad("nesting too deep") end
    skip_ws()
    local c = byte(s, pos)
    if c == 123 then return parse_object(depth)
    elseif c == 91 then return parse_array(depth)
    elseif c == 34 then return parse_string()
    elseif c == nil then bad("unexpected end of input")
    elseif find(s, "^[-0-9tfn]", pos) then bad("numbers, true, false and null are not allowed")
    else bad("unexpected character") end
  end

  if sub(s, 1, 3) == "\239\187\191" then pos = 4 end -- optional BOM
  local v = parse_value(0)
  skip_ws()
  if pos <= #s then bad("trailing data") end
  return v
end

-- names in ascending order of their UTF-8 bytes, whatever the locale
local function byteless(a, b)
  for i = 1, math.min(#a, #b) do
    local x, y = byte(a, i), byte(b, i)
    if x ~= y then return x < y end
  end
  return #a < #b
end

local function quote(s) return "'" .. (s:gsub("'", "'\\''")) .. "'" end

-- RFC 4 (validation of the whole document) and RFC 5 (as shell code), in one walk.
local function walk(o, dir, out)
  local keys = {}
  for _, k in ipairs(o.keys) do keys[#keys + 1] = k end
  table.sort(keys, byteless)
  for _, name in ipairs(keys) do
    local path = dir .. "/" .. name
    if name == "" or name == "." or name == ".." or find(name, "/", 1, true) or find(name, "\0", 1, true) then
      fail(path .. ": invalid name")
    end
    local v, q = o.map[name], quote(path)
    out[#out + 1] = "_json2dir_remove " .. q
    if type(v) == "string" then
      if find(v, "\0", 1, true) then fail(path .. ": content with NUL cannot be passed to the shell") end
      out[#out + 1] = "_json2dir_file " .. q .. " " .. quote(v)
    elseif v.arr then
      if #v ~= 2 or type(v[1]) ~= "string" or type(v[2]) ~= "string" then
        fail(path .. ': an array must be ["link", target] or ["script", content]')
      end
      if v[1] ~= "link" and v[1] ~= "script" then fail(path .. ": unknown array kind " .. v[1]) end
      if find(v[2], "\0", 1, true) then fail(path .. ": " .. v[1] .. " payload with NUL cannot be passed to the shell") end
      out[#out + 1] = "_json2dir_" .. v[1] .. " " .. q .. " " .. quote(v[2])
    else
      out[#out + 1] = "_json2dir_directory " .. q
      walk(v, path, out)
    end
  end
end

-- Only constant code here; every name, content and target is a single-quoted word.
local helpers = [[
set -eu
_json2dir_error() {
  printf 'Error: couldn\047t create %s at %s.\n' "$1" "$2" >&2
  exit 1
}
_json2dir_remove() {
  rm -f -- "$1" 2>/dev/null || :
}
_json2dir_directory() {
  if [ ! -d "$1" ]; then
    mkdir -- "$1" || _json2dir_error directory "$1"
  fi
  (cd -- "$1") || _json2dir_error directory "$1"
}
_json2dir_file() {
  printf '%s' "$2" > "$1" || _json2dir_error file "$1"
}
_json2dir_script() {
  printf '%s' "$2" > "$1" || _json2dir_error script "$1"
  chmod a+x "$1" || _json2dir_error script "$1"
}
_json2dir_link() {
  if [ -d "$1" ] && [ ! -L "$1" ]; then
    _json2dir_error symlink "$1"
  fi
  ln -s -- "$2" "$1" || _json2dir_error symlink "$1"
}]]

json2dir = {}

-- Called from TeX: read, check and translate the input; on any error report and exit 1.
function json2dir.run(file)
  local ok, err = pcall(function()
    -- --safer disables io.open entirely, but io.input(name) still opens a file for
    -- reading (fopen "r": on POSIX the bytes come through unchanged).
    local okf, f = pcall(io.input, file)
    if not okf or not f then fail("couldn't read the input file") end
    local s = f:read("a")
    f:close()
    if not s then fail("couldn't read the input file") end
    check_utf8(s)
    local root = parse(s)
    if type(root) ~= "table" or not root.obj then fail("the root must be an object") end
    local out = { helpers }
    walk(root, ".", out)
    json2dir.script = table.concat(out, "\n") -- \write adds the final newline
  end)
  if not ok then
    io.stderr:write("json2dir: " .. (type(err) == "table" and err.msg or tostring(err)) .. "\n")
    io.stderr:flush()
    os.exit(1)
  end
end
}@endgroup
%
\newlinechar=10
\directlua{json2dir.run(token.get_macro("inputfile"))}
\immediate\openout1=json2dir-tree.sh
\immediate\write1{\directlua{tex.sprint(-2, json2dir.script)}}
\immediate\closeout1
\end
'
% It reads \inputfile, validates the whole document and writes DIR/json2dir-tree.sh, a
% POSIX shell script that creates the tree. Any error: message on stderr, exit 1, no script.
%
% TeX drives; the JSON work is done by LuaTeX's embedded Lua in \directlua. --safer
% disables io.open (it uses io.input to read), so the script is written with \openout/\write: LuaTeX
% writes characters to \write files as raw UTF-8 bytes (no ^^ notation, no line wrapping),
% \newlinechar=10 turns LF into LF, and tex.sprint(-2, ...) turns every byte into a plain
% "other" character token, so the bytes come out exactly as generated.
\catcode`\{=1 \catcode`\}=2
%
% Read the Lua code with @ as the escape character and \ % and LF as ordinary characters,
% so backslashes, Lua patterns and "--" comments survive. Braces must stay balanced.
\begingroup \catcode`\%=12 \catcode10=12 \endlinechar=10 \catcode`\@=0 \catcode`\\=12
@directlua{
local byte, sub, char, find = string.byte, string.sub, string.char, string.find

local function fail(msg) error({ msg = msg }, 0) end

-- RFC 3: strict UTF-8
local function check_utf8(s)
  local i, n = 1, #s
  while i <= n do
    local c = byte(s, i)
    if c < 0x80 then
      i = i + 1
    else
      local len, min
      if c >= 0xC2 and c <= 0xDF then len, min = 2, 0x80
      elseif c >= 0xE0 and c <= 0xEF then len, min = 3, 0x800
      elseif c >= 0xF0 and c <= 0xF4 then len, min = 4, 0x10000
      else fail("input is not valid UTF-8") end
      if i + len - 1 > n then fail("input is not valid UTF-8") end
      local cp = c & (0x7F >> len)
      for k = 1, len - 1 do
        local d = byte(s, i + k)
        if d < 0x80 or d > 0xBF then fail("input is not valid UTF-8") end
        cp = cp * 64 + (d - 0x80)
      end
      if cp < min or cp > 0x10FFFF or (cp >= 0xD800 and cp <= 0xDFFF) then
        fail("input is not valid UTF-8")
      end
      i = i + len
    end
  end
end

-- RFC 3, 4: strict JSON. Objects -> {obj=true, keys=..., map=...}, arrays -> {arr=true, ...},
-- strings -> Lua strings. Numbers, true, false, null are rejected right away (RFC 4.5).
local function parse(s)
  local pos = 1
  local escapes = { [34] = '"', [92] = "\\", [47] = "/", [98] = "\b", [102] = "\f",
    [110] = "\n", [114] = "\r", [116] = "\t" }

  local function skip_ws() pos = find(s, "[^ \t\r\n]", pos) or #s + 1 end
  local function bad(what) fail("input is not valid JSON: " .. what .. " at byte " .. pos) end

  local function hex4()
    local h = sub(s, pos, pos + 3)
    if not h:match("^%x%x%x%x$") then bad("bad \\u escape") end
    pos = pos + 4
    return tonumber(h, 16)
  end

  local function parse_string()
    pos = pos + 1
    local buf = {}
    while true do
      local j = find(s, '["\\%c]', pos)
      if not j then bad("unterminated string") end
      if j > pos then buf[#buf + 1] = sub(s, pos, j - 1) end
      local c = byte(s, j)
      pos = j + 1
      if c == 34 then
        return table.concat(buf)
      elseif c == 92 then
        local e = byte(s, pos)
        if e == 117 then
          pos = pos + 1
          local cp = hex4()
          if cp >= 0xD800 and cp <= 0xDBFF then
            if sub(s, pos, pos + 1) ~= "\\u" then bad("unpaired surrogate") end
            pos = pos + 2
            local lo = hex4()
            if lo < 0xDC00 or lo > 0xDFFF then bad("unpaired surrogate") end
            cp = 0x10000 + (cp - 0xD800) * 0x400 + (lo - 0xDC00)
          elseif cp >= 0xDC00 and cp <= 0xDFFF then
            bad("unpaired surrogate")
          end
          buf[#buf + 1] = utf8.char(cp)
        elseif e and escapes[e] then
          buf[#buf + 1] = escapes[e]
          pos = pos + 1
        else
          bad("bad escape")
        end
      elseif c == 127 then
        buf[#buf + 1] = "\127"
      else
        bad("control character in string")
      end
    end
  end

  local parse_value

  local function parse_object(depth)
    pos = pos + 1
    local o = { obj = true, keys = {}, map = {} }
    skip_ws()
    if byte(s, pos) == 125 then pos = pos + 1; return o end
    while true do
      skip_ws()
      if byte(s, pos) ~= 34 then bad("expected member name") end
      local name = parse_string()
      skip_ws()
      if byte(s, pos) ~= 58 then bad("expected ':'") end
      pos = pos + 1
      local v = parse_value(depth + 1)
      if o.map[name] == nil then o.keys[#o.keys + 1] = name end
      o.map[name] = v -- duplicate names: the last occurrence wins (RFC 4.2)
      skip_ws()
      local c = byte(s, pos)
      pos = pos + 1
      if c == 125 then return o end
      if c ~= 44 then pos = pos - 1; bad("expected ',' or end of object") end
    end
  end

  local function parse_array(depth)
    pos = pos + 1
    local a = { arr = true }
    skip_ws()
    if byte(s, pos) == 93 then pos = pos + 1; return a end
    while true do
      a[#a + 1] = parse_value(depth + 1)
      skip_ws()
      local c = byte(s, pos)
      pos = pos + 1
      if c == 93 then return a end
      if c ~= 44 then pos = pos - 1; bad("expected ',' or ']'") end
    end
  end

  function parse_value(depth)
    if depth > 1000 then bad("nesting too deep") end
    skip_ws()
    local c = byte(s, pos)
    if c == 123 then return parse_object(depth)
    elseif c == 91 then return parse_array(depth)
    elseif c == 34 then return parse_string()
    elseif c == nil then bad("unexpected end of input")
    elseif find(s, "^[-0-9tfn]", pos) then bad("numbers, true, false and null are not allowed")
    else bad("unexpected character") end
  end

  if sub(s, 1, 3) == "\239\187\191" then pos = 4 end -- optional BOM
  local v = parse_value(0)
  skip_ws()
  if pos <= #s then bad("trailing data") end
  return v
end

-- names in ascending order of their UTF-8 bytes, whatever the locale
local function byteless(a, b)
  for i = 1, math.min(#a, #b) do
    local x, y = byte(a, i), byte(b, i)
    if x ~= y then return x < y end
  end
  return #a < #b
end

local function quote(s) return "'" .. (s:gsub("'", "'\\''")) .. "'" end

-- RFC 4 (validation of the whole document) and RFC 5 (as shell code), in one walk.
local function walk(o, dir, out)
  local keys = {}
  for _, k in ipairs(o.keys) do keys[#keys + 1] = k end
  table.sort(keys, byteless)
  for _, name in ipairs(keys) do
    local path = dir .. "/" .. name
    if name == "" or name == "." or name == ".." or find(name, "/", 1, true) or find(name, "\0", 1, true) then
      fail(path .. ": invalid name")
    end
    local v, q = o.map[name], quote(path)
    out[#out + 1] = "_json2dir_remove " .. q
    if type(v) == "string" then
      if find(v, "\0", 1, true) then fail(path .. ": content with NUL cannot be passed to the shell") end
      out[#out + 1] = "_json2dir_file " .. q .. " " .. quote(v)
    elseif v.arr then
      if #v ~= 2 or type(v[1]) ~= "string" or type(v[2]) ~= "string" then
        fail(path .. ': an array must be ["link", target] or ["script", content]')
      end
      if v[1] ~= "link" and v[1] ~= "script" then fail(path .. ": unknown array kind " .. v[1]) end
      if find(v[2], "\0", 1, true) then fail(path .. ": " .. v[1] .. " payload with NUL cannot be passed to the shell") end
      out[#out + 1] = "_json2dir_" .. v[1] .. " " .. q .. " " .. quote(v[2])
    else
      out[#out + 1] = "_json2dir_directory " .. q
      walk(v, path, out)
    end
  end
end

-- Only constant code here; every name, content and target is a single-quoted word.
local helpers = [[
set -eu
_json2dir_error() {
  printf 'Error: couldn\047t create %s at %s.\n' "$1" "$2" >&2
  exit 1
}
_json2dir_remove() {
  rm -f -- "$1" 2>/dev/null || :
}
_json2dir_directory() {
  if [ ! -d "$1" ]; then
    mkdir -- "$1" || _json2dir_error directory "$1"
  fi
  (cd -- "$1") || _json2dir_error directory "$1"
}
_json2dir_file() {
  printf '%s' "$2" > "$1" || _json2dir_error file "$1"
}
_json2dir_script() {
  printf '%s' "$2" > "$1" || _json2dir_error script "$1"
  chmod a+x "$1" || _json2dir_error script "$1"
}
_json2dir_link() {
  if [ -d "$1" ] && [ ! -L "$1" ]; then
    _json2dir_error symlink "$1"
  fi
  ln -s -- "$2" "$1" || _json2dir_error symlink "$1"
}]]

json2dir = {}

-- Called from TeX: read, check and translate the input; on any error report and exit 1.
function json2dir.run(file)
  local ok, err = pcall(function()
    -- --safer disables io.open entirely, but io.input(name) still opens a file for
    -- reading (fopen "r": on POSIX the bytes come through unchanged).
    local okf, f = pcall(io.input, file)
    if not okf or not f then fail("couldn't read the input file") end
    local s = f:read("a")
    f:close()
    if not s then fail("couldn't read the input file") end
    check_utf8(s)
    local root = parse(s)
    if type(root) ~= "table" or not root.obj then fail("the root must be an object") end
    local out = { helpers }
    walk(root, ".", out)
    json2dir.script = table.concat(out, "\n") -- \write adds the final newline
  end)
  if not ok then
    io.stderr:write("json2dir: " .. (type(err) == "table" and err.msg or tostring(err)) .. "\n")
    io.stderr:flush()
    os.exit(1)
  end
end
}@endgroup
%
\newlinechar=10
\directlua{json2dir.run(token.get_macro("inputfile"))}
\immediate\openout1=json2dir-tree.sh
\immediate\write1{\directlua{tex.sprint(-2, json2dir.script)}}
\immediate\closeout1
\end
